\section{Follow-up checks and analyses}\label{sec:review}
This note reports follow-up checks and analyses, all specified after every test had been scored. The analyses were registered together, with their code and data, before any of their quantities was computed. None changes a published estimate. Two runs departed from the planned order of execution without any change to code, setting or value.

\subsection*{Scoring of repeated draws}
At each budget a test draws paired cells several times. For predictions $C_1,\dots,C_D$ from $D$ draws and held-out truth $T$, the numerator of the recovered fraction is $G(C)=2\langle C,T\rangle-\|C\|^2$, and
\[
G\Big(\frac1D\sum_dC_d\Big)=\frac1D\sum_dG(C_d)+\frac1D\sum_d\Big\|C_d-\frac1D\sum_{d'}C_{d'}\Big\|^2 ,
\]
so scoring the averaged prediction would credit a study with the paired cells of all its draws. Each curve in this paper is the mean of per-draw scores: the tests store each arm's mean prediction and the mean over draws of $\|C_d\|^2$, from which $2\langle\bar C,T\rangle-\overline{\|C\|^2}=\frac1D\sum_dG(C_d)$. A check reruns the registered prediction code of each test unchanged, observes every draw's prediction as the code accumulates it, scores each on its own and compares (Supplementary Table~\ref{tab:draws}). Across \RvDrawTests{} tests and data sets, the curves averaged over per-draw scores equalled the published curves to within \RvDrawMaxDiff, the regenerated means and squared norms equalled the stored ones, and the rerun of the bone-marrow test reproduced its recorded predictions exactly. Scoring the averaged predictions instead would have raised the curves by a median of \RvDrawInflMedianMin{} to \RvDrawInflMedianMax{} percentage points per test, and by up to \RvDrawInflMax{} points where draws disagreed most. The development comparison and the validation's readout score each draw inside the draw loop. A second check recomputes every curve from the stored per-draw scores without the data.

\begin{table}[h]
\caption{\textbf{Every recovery curve is the mean of per-draw scores.} Per test or data set: held-out units whose predictions were regenerated (of all units; each benchmark data set was rerun for the first of its three folds), draws per budget, the largest absolute difference between the curve averaged over per-draw scores and the published curve (or, for partial reruns, the fold's stored terms), the largest relative difference between the stored and regenerated mean squared norms, and the median and largest amount, in percentage points of recovered fraction, by which scoring the averaged prediction would have exceeded the published curve. External test: every draw's prediction is stored.}
\label{tab:draws}
\centering\footnotesize
% Generated by extension/synthesis/make_assets.py. Do not edit.
\begin{tabular}{@{}lrrrrrr@{}}
\toprule
 & Units & Draws & \multicolumn{2}{c}{Largest difference} & \multicolumn{2}{c}{Averaged prediction} \\
\cmidrule(lr){4-5}\cmidrule(l){6-7}
Test & rerun & & curve & msq & median & max \\
\midrule
External test & 1/1 & 20 & $3\times10^{-16}$ & -- & 9.0 & 61 \\
Bone-marrow test & 12/12 & 5 & $3\times10^{-9}$ & $4\times10^{-16}$ & 28.9 & 537 \\
Validation test & 34/34 & 5 & $4\times10^{-9}$ & $10^{-16}$ & 9.1 & 360 \\
Benchmark: banc & 1/3 & 5 & $4\times10^{-16}$ & $3\times10^{-16}$ & 8.0 & 32 \\
Benchmark: banc\_crossanimal & 1/3 & 5 & $6\times10^{-16}$ & $4\times10^{-16}$ & 27.5 & 123 \\
Benchmark: bmmc\_cite & 1/3 & 5 & $2\times10^{-15}$ & $3\times10^{-16}$ & 8.0 & 86 \\
Benchmark: bmmc\_multiome & 1/3 & 5 & $2\times10^{-16}$ & $4\times10^{-16}$ & 2.3 & 43 \\
Benchmark: colon & 1/3 & 5 & $10^{-16}$ & $4\times10^{-16}$ & 4.3 & 31 \\
Benchmark: gouwens\_visp & 1/3 & 5 & $10^{-16}$ & $3\times10^{-16}$ & 6.2 & 16 \\
Benchmark: hao & 1/3 & 5 & $6\times10^{-16}$ & $4\times10^{-16}$ & 7.1 & 208 \\
Benchmark: malecns & 1/3 & 5 & $4\times10^{-16}$ & $3\times10^{-16}$ & 9.5 & 33 \\
Benchmark: scala\_m1 & 1/3 & 5 & $10^{-16}$ & $4\times10^{-16}$ & 5.5 & 29 \\
Benchmark: stephenson & 1/3 & 5 & $2\times10^{-16}$ & $4\times10^{-16}$ & 5.1 & 39 \\
Law: banc\_vpn & 2 problems & None & $4\times10^{-16}$ & $3\times10^{-16}$ & 11.3 & 19 \\
Law: bmcite & 3 problems & None & $8\times10^{-16}$ & $4\times10^{-16}$ & 2.7 & 25 \\
Law: fafb\_vpn & 3 problems & None & $6\times10^{-16}$ & $4\times10^{-16}$ & 13.6 & 25 \\
Law: fetal\_cortex & 3 problems & None & $10^{-16}$ & $4\times10^{-16}$ & 4.7 & 51 \\
Law: human\_gaba & 3 problems & None & $2\times10^{-16}$ & $4\times10^{-16}$ & 19.5 & 50 \\
Law: malt10k & 3 problems & None & $10\times10^{-16}$ & $5\times10^{-16}$ & 3.4 & 24 \\
Law: pbmc10k & 3 problems & None & $10^{-15}$ & $4\times10^{-16}$ & 4.1 & 27 \\
Law: sln111 & 3 problems & None & $4\times10^{-16}$ & $6\times10^{-16}$ & 10.3 & 67 \\
Law: sln206 & 3 problems & None & $6\times10^{-16}$ & $6\times10^{-16}$ & 8.5 & 33 \\
Law: snare\_cortex & 3 problems & None & $7\times10^{-16}$ & $4\times10^{-16}$ & 5.0 & 82 \\
\bottomrule
\end{tabular}

\end{table}

\subsection*{The noise estimate after contrasts}
Proposition~\ref{prop:oracle} assumes independent product rows, as when paired cells are centred on means estimated from other cells (development, the external test, the benchmark and the prospective test of the law). The validation test and the reanalyses of the bone-marrow and fly tests centre each population by Helmert contrasts. Contrast rows are uncorrelated, but their products are dependent unless the cells are Gaussian, so the noise estimate $t_b$ of Proposition~\ref{prop:oracle} need not be unbiased for them. Its exact expectation is as follows.

Consider one block and populations $p=1,\dots,P$ with $n_p\ge2$ independent cells each, the scaling and the bases held fixed. Population $p$ contributes $m_p=n_p-1$ rows $h_{pj}=\sum_iH^{(p)}_{ji}u_{pi}$, where $H^{(p)}$ has orthonormal rows orthogonal to the vector of ones, and $z_{pj}$ is the block of $h^x_{pj}h^{y\top}_{pj}$. With $N=\sum_pm_p$ rows, $m_b$ is their mean and $t_b=\sum\|z_{pj}-m_b\|^2/\{N(N-1)\}$. For a coefficient $(k,l)$ of the block, let $\sigma^{(p)}_{kl}$ be the population's covariance of $\tilde x_k$ and $\tilde y_l$ and $\kappa^{(p)}_{kl}=\mathbb{E}[\bar x_k^2\bar y_l^2]-\operatorname{Var}\tilde x_k\operatorname{Var}\tilde y_l-2(\sigma^{(p)}_{kl})^2$ their fourth cumulant ($\bar x,\bar y$ centred at the population means). Write $\mu_p$ for the block of $\sigma^{(p)}$, $\bar\mu=\sum_pm_p\mu_p/N$, $\gamma_p=\sum_{(k,l)}\{\operatorname{Var}\tilde x_k\operatorname{Var}\tilde y_l+(\sigma^{(p)}_{kl})^2\}$, $\kappa_p=\sum_{(k,l)}\kappa^{(p)}_{kl}$ and $F_p=\sum_{j,i}(H^{(p)}_{ji})^4$.

\begin{proposition}[The noise estimate after orthonormal contrasts]\label{prop:contrast}
\[
\operatorname{tr}\operatorname{Cov}(m_b)=\frac{1}{N^2}\sum_p\Big(m_p\gamma_p+\frac{m_p^2}{n_p}\kappa_p\Big),\qquad
\mathbb{E}t_b-\operatorname{tr}\operatorname{Cov}(m_b)=\frac{\sum_pm_p\|\mu_p-\bar\mu\|^2-\sum_p\kappa_p\big(m_p^2/n_p-F_p\big)}{N(N-1)} .
\]
For Helmert contrasts $F_p=m_p+3-2\mathcal H_{m_p}-3/n_p$, with $\mathcal H_m$ the $m$-th harmonic number, so that $m_p^2/n_p-F_p=2\mathcal H_{m_p}-4m_p/n_p$, which lies between 0 and $2\ln m_p$.
\end{proposition}
\noindent\textit{Proof.} Contrasts remove the population mean, so take the cells of population $p$ centred, with coordinates $a=\tilde x_k$ and $b=\tilde y_l$ for one coefficient. For independent centred cells, $\mathbb{E}[a_ib_{i'}a_{r}b_{r'}]=\sigma_{ab}^2\delta_{ii'}\delta_{rr'}+\sigma_{aa}\sigma_{bb}\delta_{ir}\delta_{i'r'}+\sigma_{ab}^2\delta_{ir'}\delta_{i'r}+\kappa_{ab}\delta_{ii'rr'}$, where the last term is one when all four indices coincide. Summing against $H_{ji}H_{ji'}H_{kr}H_{kr'}$ and using $\sum_iH_{ji}H_{ki}=\delta_{jk}$ gives $\mathbb{E}z_j=\sigma_{ab}$ and $\operatorname{Cov}(z_j,z_k)=\delta_{jk}(\sigma_{aa}\sigma_{bb}+\sigma_{ab}^2)+\kappa_{ab}\sum_iH_{ji}^2H_{ki}^2$. Rows of different populations are independent. Summing over the coefficients of the block, over $j$ and over $k$, with $\sum_jH_{ji}^2=1-1/n_p$, gives the variance of $m_b$. For rows with means $\mu_j$, $\mathbb{E}\sum_j\|z_j-m_b\|^2=\sum_j\operatorname{tr}\operatorname{Cov}z_j+\sum_j\|\mu_j-\bar\mu\|^2-N\operatorname{tr}\operatorname{Cov}(m_b)$, and dividing by $N(N-1)$ and subtracting leaves the covariances between distinct rows, $\sum_{j\ne k}\sum_iH_{ji}^2H_{ki}^2=m_p^2/n_p-F_p$ within population $p$. For the Helmert row $j$, with $j$ entries $1/\sqrt{j(j+1)}$ and one entry $-j/\sqrt{j(j+1)}$, $\sum_iH_{ji}^4=1+1/j-3/(j+1)$, and summing over $j$ gives $F_p$. Finally $2\mathcal H_m-4m/(m+1)$ is 0 at $m=1$, increases with $m$ because $2/(m+1)\ge4/\{(m+1)(m+2)\}$, and is at most $2\ln m$ because $\mathcal H_m\le1+\ln m$. $\square$

The first term of the bias is non-negative: populations whose cross-covariances differ inflate the estimate, and shrinkage becomes more cautious. The second vanishes for Gaussian cells ($\kappa_p=0$) and otherwise has the sign opposite to the fourth cumulant. Relative to the fourth-cumulant part of $\operatorname{tr}\operatorname{Cov}(m_b)$, it is of order $(\ln m_p)/m_p$ for Helmert contrasts, whose rows are dominated by a single new cell; for a basis with entries of equal size ($F_p=m_p/n_p$) it would miss that part entirely. The referee's example, a population of three cells with $x\in\{-3,0,3\}$ of probabilities $1/18,8/9,1/18$ and $y=x\pm1$, has $\gamma=3$, $\kappa=6$, $\operatorname{Var}(m)=3/2+6/3=3.5$ and $\mathbb{E}t=3.5-6\cdot(1/3)/2=2.5$. The proposition holds the scaling fixed, whereas the implemented pipeline estimates the standard deviations from the same contrasts; the calibration below includes that step.

\paragraph{Calibration.} For every fold of the validation test (atlas axes) and of the bone-marrow test (other-site axes, the reanalysis with contrasts) and each budget, 500 independent resamples of paired cells were drawn with replacement from the fold's reservoir and passed through the complete pipeline: Helmert contrasts in draw order, scaling by the contrasts' standard deviations, rotation and blocks. The calibration ratio is the mean of a noise estimate over resamples divided by the variance of the block mean across resamples, pooled over folds, with 95\% intervals from resampling donors (validation) or held-out batches (bone marrow) and then resamples within folds (Supplementary Tables~\ref{tab:noisecal} and~\ref{tab:noiseblocks}). The implemented estimate was slightly conservative at every budget: its ratio, summed over blocks, was \RvCalValMin{} to \RvCalValMax{} in the validation test and \RvCalBmMin{} to \RvCalBmMax{} in the bone-marrow test, and \RvCalValBlockMin{} to \RvCalValBlockMax{} and \RvCalBmBlockMin{} to \RvCalBmBlockMax{} block by block, highest in the blocks of the leading RNA axes, where the populations' cross-covariances differ most, as the first term of Proposition~\ref{prop:contrast} predicts; the excess grew with the budget. Random orthonormal contrasts gave \RvCalValOrthoMin{} to \RvCalValOrthoMax. Reordering the cells within populations left the ratios unchanged, as it must, since the order changes the estimate's variance but not its expectation; the variance was appreciable at small budgets, where the most affected block's estimate varied over orders with a median coefficient of variation of \RvCalValOrderCvMax\% at 25 paired cells, \RvCalValOrderCvHundred\% at 100 and \RvCalValOrderCvEight\% at 800, and its shrinkage factor by a median range of \RvCalValOrderFactorMax, \RvCalValOrderFactorHundred{} and \RvCalValOrderFactorEight. The jackknife over cells, a dependence-aware alternative, overestimated the noise at small budgets (ratio \RvCalValJackMax{} at 25 paired cells, falling to \RvCalValJackMin{} at 800), and centring on the drawn cells' own averages underestimated it (\RvCalValOwnMin{} at 25), reproducing the failure of the bone-marrow test's original centring. The jackknife's larger estimates shrink more, and on the same resamples they raised the recovered fraction by \RvCalValJackGainMin{} to \RvCalValJackGainMax{} percentage points in the validation test and by \RvCalBmJackGainMin{} to \RvCalBmJackGainMax{} in the bone-marrow test, most at small budgets (Supplementary Table~\ref{tab:noisecal}): at these budgets the positive-part rule with an unbiased noise estimate shrinks less than would be best. The implemented estimate, fixed before the tests, is kept in every result reported here; its calibration after contrasts is empirical, and Proposition~\ref{prop:contrast} gives the form and order of its bias.

\begin{table}[h]
\caption{\textbf{Calibration of the noise estimate after contrasts.} Mean noise estimate over 500 independent resamples of paired cells, divided by the variance of the block mean across the same resamples, summed over the eight blocks and pooled over folds (unrounded to three decimals). For the implemented estimate (Helmert contrasts in draw order): its 95\% interval from resampling donors (validation) or held-out batches (bone marrow) and resamples within folds; the range of its ratio over the eight blocks; and, for the block most affected by the order of cells, the median coefficient of variation of the estimate and the median range of the block's shrinkage factor over 20 random orders of the cells within populations. Alternatives on the same resamples: Helmert contrasts in a random order, random orthonormal contrasts, the jackknife over cells, and centring on the drawn cells' own averages ($\sqrt{n/(n-1)}$ times the deviations, as in the bone-marrow test, compared with the variance of its own block mean). Last rows: fraction of the dependence recovered by the atlas estimate with condition maps (validation) or the other-site estimate (bone marrow) on the same resamples, with the implemented and with the jackknife noise.}
\label{tab:noisecal}
\centering\scriptsize
% Generated by extension/synthesis/make_assets.py. Do not edit.
\begin{tabular}{@{}lrrrrrr@{}}
\toprule
 & \multicolumn{6}{c}{Paired cells} \\
\cmidrule(l){2-7}
Noise estimate & 25 & 50 & 100 & 200 & 400 & 800 \\
\midrule
\multicolumn{7}{@{}l}{\textit{Validation test}} \\
Helmert, draw order (implemented) & 1.005 & 1.008 & 1.011 & 1.014 & 1.016 & 1.016 \\
\quad 95\% interval & 1.003--1.006 & 1.006--1.010 & 1.009--1.012 & 1.012--1.015 & 1.014--1.017 & 1.015--1.018 \\
\quad range over blocks & 1.00--1.05 & 1.00--1.05 & 1.00--1.06 & 1.01--1.06 & 1.01--1.07 & 1.01--1.07 \\
\quad order, CV (largest) & 12.3\% & 10.8\% & 8.5\% & 6.2\% & 4.2\% & 2.8\% \\
\quad order, factor range & 0.223 & 0.184 & 0.103 & 0.057 & 0.032 & 0.016 \\
Helmert, random order & 1.005 & 1.008 & 1.011 & 1.014 & 1.015 & 1.016 \\
Random orthonormal & 1.003 & 1.005 & 1.007 & 1.008 & 1.010 & 1.010 \\
Jackknife over cells & 1.478 & 1.258 & 1.134 & 1.073 & 1.045 & 1.031 \\
Own-mean centring & 0.674 & 0.798 & 0.898 & 0.958 & 0.988 & 1.003 \\
Recovered (\%), implemented & \ensuremath{-}1.3 & 8.8 & 19 & 28 & 33 & 35 \\
Recovered (\%), jackknife & 1.3 & 10 & 22 & 31 & 34 & 35 \\
\addlinespace
\multicolumn{7}{@{}l}{\textit{Bone-marrow test}} \\
Helmert, draw order (implemented) & 1.005 & 1.012 & 1.016 & 1.024 & 1.029 & 1.029 \\
\quad 95\% interval & 1.003--1.008 & 1.008--1.016 & 1.012--1.021 & 1.018--1.030 & 1.021--1.036 & 1.022--1.036 \\
\quad range over blocks & 1.00--1.08 & 1.00--1.16 & 1.00--1.18 & 1.01--1.26 & 1.01--1.28 & 1.01--1.31 \\
\quad order, CV (largest) & 12.4\% & 11.6\% & 9.7\% & 7.8\% & 6.0\% & 4.4\% \\
\quad order, factor range & 0.218 & 0.148 & 0.113 & 0.068 & 0.032 & 0.015 \\
Helmert, random order & 1.005 & 1.012 & 1.016 & 1.024 & 1.028 & 1.029 \\
Random orthonormal & 1.005 & 1.010 & 1.015 & 1.023 & 1.030 & 1.034 \\
Jackknife over cells & 1.538 & 1.328 & 1.193 & 1.117 & 1.075 & 1.051 \\
Own-mean centring & 0.650 & 0.759 & 0.855 & 0.934 & 0.982 & 1.006 \\
Recovered (\%), implemented & \ensuremath{-}4.2 & 3.8 & 10 & 15 & 18 & 17 \\
Recovered (\%), jackknife & 1.7 & 7.8 & 15 & 20 & 20 & 18 \\
\bottomrule
\end{tabular}

\end{table}

\begin{table}[h]
\caption{\textbf{Calibration of the implemented noise estimate block by block.} As Supplementary Table~\ref{tab:noisecal}, for Helmert contrasts in draw order, in each of the eight blocks (RNA axes 1--5, 6--20, 21--60 and the rest, crossed with protein axes 1--3 and the rest), with 95\% intervals.}
\label{tab:noiseblocks}
\centering\scriptsize
% Generated by extension/synthesis/make_assets.py. Do not edit.
\begin{tabular}{@{}lrrrrrr@{}}
\toprule
 & \multicolumn{6}{c}{Paired cells} \\
\cmidrule(l){2-7}
Block (RNA/protein axes) & 25 & 50 & 100 & 200 & 400 & 800 \\
\midrule
\multicolumn{7}{@{}l}{\textit{Validation test}} \\
1--5/1--3 & 1.048 & 1.048 & 1.050 & 1.052 & 1.061 & 1.054 \\
 & (1.032--1.062) & (1.030--1.065) & (1.030--1.070) & (1.036--1.066) & (1.047--1.076) & (1.040--1.068) \\
1--5/rest & 1.038 & 1.051 & 1.059 & 1.063 & 1.066 & 1.069 \\
 & (1.030--1.045) & (1.043--1.059) & (1.052--1.065) & (1.056--1.069) & (1.061--1.071) & (1.064--1.074) \\
6--20/1--3 & 1.011 & 1.009 & 1.015 & 1.022 & 1.025 & 1.030 \\
 & (1.004--1.018) & (1.003--1.015) & (1.009--1.022) & (1.016--1.029) & (1.019--1.031) & (1.024--1.037) \\
6--20/rest & 1.003 & 1.007 & 1.011 & 1.016 & 1.019 & 1.019 \\
 & (1.001--1.005) & (1.005--1.009) & (1.009--1.013) & (1.013--1.018) & (1.016--1.021) & (1.018--1.021) \\
21--60/1--3 & 1.005 & 1.005 & 1.010 & 1.013 & 1.015 & 1.018 \\
 & (1.002--1.009) & (1.001--1.008) & (1.006--1.014) & (1.008--1.017) & (1.011--1.019) & (1.014--1.023) \\
21--60/rest & 1.001 & 1.003 & 1.005 & 1.008 & 1.010 & 1.012 \\
 & (1.000--1.002) & (1.002--1.004) & (1.004--1.007) & (1.007--1.009) & (1.009--1.012) & (1.011--1.014) \\
rest/1--3 & 1.003 & 1.005 & 1.008 & 1.010 & 1.012 & 1.012 \\
 & (1.000--1.005) & (1.003--1.008) & (1.005--1.010) & (1.008--1.013) & (1.009--1.015) & (1.009--1.015) \\
rest/rest & 1.000 & 1.002 & 1.004 & 1.006 & 1.007 & 1.008 \\
 & (1.000--1.001) & (1.001--1.003) & (1.003--1.005) & (1.005--1.007) & (1.006--1.009) & (1.007--1.009) \\
\addlinespace
\multicolumn{7}{@{}l}{\textit{Bone-marrow test}} \\
1--5/1--3 & 1.082 & 1.161 & 1.184 & 1.265 & 1.284 & 1.310 \\
 & (1.041--1.125) & (1.116--1.206) & (1.142--1.224) & (1.203--1.321) & (1.214--1.346) & (1.236--1.375) \\
1--5/rest & 1.017 & 1.034 & 1.040 & 1.059 & 1.071 & 1.074 \\
 & (1.012--1.023) & (1.023--1.046) & (1.030--1.051) & (1.045--1.076) & (1.053--1.090) & (1.056--1.092) \\
6--20/1--3 & 1.031 & 1.049 & 1.077 & 1.099 & 1.113 & 1.113 \\
 & (1.008--1.053) & (1.031--1.069) & (1.053--1.102) & (1.070--1.128) & (1.084--1.139) & (1.087--1.136) \\
6--20/rest & 1.002 & 1.009 & 1.015 & 1.023 & 1.032 & 1.033 \\
 & (1.000--1.005) & (1.006--1.013) & (1.010--1.021) & (1.018--1.029) & (1.023--1.041) & (1.026--1.040) \\
21--60/1--3 & 1.018 & 1.028 & 1.034 & 1.038 & 1.046 & 1.042 \\
 & (1.008--1.028) & (1.019--1.036) & (1.023--1.045) & (1.027--1.048) & (1.034--1.057) & (1.030--1.054) \\
21--60/rest & 1.000 & 1.003 & 1.005 & 1.010 & 1.013 & 1.012 \\
 & (0.998--1.001) & (1.001--1.005) & (1.003--1.007) & (1.007--1.012) & (1.010--1.016) & (1.010--1.015) \\
rest/1--3 & 1.015 & 1.020 & 1.026 & 1.028 & 1.031 & 1.028 \\
 & (1.009--1.022) & (1.013--1.028) & (1.020--1.032) & (1.023--1.034) & (1.024--1.038) & (1.022--1.035) \\
rest/rest & 1.001 & 1.003 & 1.005 & 1.008 & 1.010 & 1.009 \\
 & (0.999--1.002) & (1.001--1.004) & (1.003--1.006) & (1.006--1.010) & (1.008--1.012) & (1.007--1.011) \\
\bottomrule
\end{tabular}

\end{table}

\subsection*{Axes and condition maps}
The estimator takes two things from a reference: axes, along which the paired estimate is shrunk block by block, and each cell type's covariance, through which the condition map turns the estimate pooled over types into one for each type. In the validation test, every draw of paired cells was regenerated with the test's registered code (the regenerated means and squared norms equalled the stored ones to within \RvValMeanDiff{} and \RvValMsqDiff) and scored with four versions for the atlas: axes and maps (the test's arm); axes only, the same pooled estimate given to every cell type; maps only, paired-only single-block James--Stein pooled over types and passed through the atlas's maps; and neither, paired-only single-block James--Stein pooled over types. Axes only was also scored for the other-pool and own-sample references.

Each alone did little (Supplementary Table~\ref{tab:axesmaps}). With 100 paired cells, the axes recovered \RvAxesHundred\% of the dependence within cell types (95\% CI, \RvAxesHundredLo--\RvAxesHundredHi\%) and the maps \RvMapOnlyHundred\% (\RvMapOnlyHundredLo--\RvMapOnlyHundredHi\%), against \RvNeitherHundred\% with neither and \RvBothHundred\% (\RvBothHundredLo--\RvBothHundredHi\%) with both. With 800, the axes recovered \RvAxesEight\%, the maps \RvMapOnlyEight\% and both \RvBothEight\%. Neither alone reached the primary target within 800 paired cells, whereas both together reached it with \RvBothCells{} (95\% CI, \RvBothCellsLo--\RvBothCellsHi; the regenerated draws add the readout's budget of 150 paired cells, which moves the interpolation slightly from the prespecified test's \VlCellsAtlas). The same held for the study's own unpaired cells and its other pools: without maps they recovered \RvOwnAxesEight\% and \RvOtherAxesEight\% with 800 paired cells. An estimate pooled over cell types can recover only the part of the within-type dependence that the types share; the maps make it type-specific, and the axes make the pooled estimate accurate from few cells.

\paragraph{Reweighting the reference.} A condition map passes the pooled paired estimate through the reference's pooled covariance, which weights cell types by the reference's cells, whereas the pooled contrast estimate weights them by their contrast rows (the sum over a type's populations of $n_p-1$). The reweighted map moves the reference's pooled latent covariance $S$ to $S+\sum_c(w_c-w^0_c)S_c$, with $S_c$ the reference's covariance of type $c$ in pooled standard deviations, $w^0_c$ the reference's share of type $c$ and $w_c$ the draw's share of contrast rows, rescales it to the reweighted pooled standard deviations and rebuilds the maps exactly as the unweighted ones are built; with $w=w^0$ it gives the unweighted maps, which the code checks (largest difference 0). It was applied to the full atlas and to the two skewed atlases of Supplementary Table~\ref{tab:robust}, in every draw. Reweighting rescued both skewed atlases: with B cells and CD14 monocytes from 15 patients, the estimator needed \RvFewBMonoRwCells{} paired cells (95\% CI, \RvFewBMonoRwCellsLo--\RvFewBMonoRwCellsHi) with reweighted maps against \RvFewBMonoMappedCells{} with the unweighted ones, and with CD4 and CD8 T cells from 15 patients \RvFewTRwCells{} (\RvFewTRwCellsLo--\RvFewTRwCellsHi) against \RvFewTMappedCells; the full atlas needed \RvFullRwCells{} with reweighted maps and \RvBothCells{} with its own. With 25 paired cells, when some cell types had almost no contrast rows, the reweighted maps were far worse than predicting no dependence for every atlas (\RvFullRwTwentyFive\% for the full atlas).

\begin{table}[h]
\caption{\textbf{Axes and condition maps in the validation test.} Fraction of the within-type dependence recovered (\%), pooled over the \VlNFolds{} held-out samples, and paired cells to the test's primary target (\VlTarget\%) with 95\% intervals from resampling donors. Top: the atlas with its axes and maps (the test's arm), its axes only, its maps only (paired-only single-block James--Stein pooled over types, through the atlas's maps) and neither; the other pools' and the own samples' unpaired cells with and without maps; the best paired-only method at each budget. Bottom: the unchanged atlas and the two atlases with a skewed mix of cell types (Supplementary Table~\ref{tab:robust}), with their unweighted maps, with maps reweighted to each draw's contrast rows per cell type, and without maps.}
\label{tab:axesmaps}
\centering\footnotesize
% Generated by extension/synthesis/make_assets.py. Do not edit.
\begin{tabular}{@{}lrrrrrrrr@{}}
\toprule
 & \multicolumn{7}{c}{Recovered (\%) with paired cells} & Paired cells \\
\cmidrule(lr){2-8}
Reference and version & 25 & 50 & 100 & 150 & 200 & 400 & 800 & to the target (95\% CI) \\
\midrule
\multicolumn{9}{@{}l}{Atlas} \\
\quad Axes and maps & 0.5 & 11 & 23 & 29 & 33 & 42 & 48 & 145 (124--183) \\
\quad Axes only & \ensuremath{-}0.7 & 4.0 & 7.8 & 9.3 & 10 & 13 & 14 & $>$800 \\
\quad Maps only & \ensuremath{-}0.2 & 1.8 & 4.7 & 7.1 & 9.5 & 18 & 28 & $>$800 \\
\quad Neither & 0.0 & 1.1 & 2.4 & 3.4 & 4.4 & 7.1 & 9.9 & $>$800 \\
\addlinespace
\multicolumn{9}{@{}l}{Other pools and donors} \\
\quad Axes and maps & 8.8 & 23 & 34 & 39 & 43 & 50 & 57 & 70 (58--85) \\
\quad Axes only & 1.6 & 7.5 & 10 & 11 & 12 & 13 & 14 & $>$800 \\
\addlinespace
\multicolumn{9}{@{}l}{Own samples} \\
\quad Axes and maps & 5.6 & 19 & 27 & 31 & 34 & 40 & 45 & 109 (82--165) \\
\quad Axes only & 0.5 & 6.5 & 9.2 & 9.7 & 10 & 12 & 13 & $>$800 \\
\addlinespace
\multicolumn{9}{@{}l}{Paired cells only} \\
\quad Best of eight methods & 0.5 & 2.4 & 3.5 & 5.1 & 7.7 & 14 & 13 & $>$800 \\
\midrule
\multicolumn{9}{@{}l}{\textit{Reweighting the atlas}} \\
\multicolumn{9}{@{}l}{All 118 atlas patients} \\
\quad Unweighted maps & 0.5 & 11 & 23 & 29 & 33 & 42 & 48 & 145 (124--183) \\
\quad Reweighted maps & \ensuremath{-}37 & 8.3 & 24 & 31 & 35 & 44 & 50 & 131 (114--162) \\
\quad No maps & \ensuremath{-}0.7 & 4.0 & 7.8 & 9.3 & 10 & 13 & 14 & $>$800 \\
\addlinespace
\multicolumn{9}{@{}l}{B cells and CD14 monocytes from 15} \\
\quad Unweighted maps & \ensuremath{-}13 & \ensuremath{-}6.1 & 2.8 & 5.7 & 10.0 & 17 & 20 & $>$800 \\
\quad Reweighted maps & \ensuremath{-}25 & 10.0 & 23 & 29 & 32 & 40 & 46 & 148 (122--192) \\
\quad No maps & \ensuremath{-}0.5 & 4.2 & 8.0 & 9.3 & 10 & 12 & 14 & $>$800 \\
\addlinespace
\multicolumn{9}{@{}l}{CD4 and CD8 T cells from 15} \\
\quad Unweighted maps & \ensuremath{-}3.3 & 7.2 & 14 & 15 & 14 & 8.8 & \ensuremath{-}2.7 & $>$800 \\
\quad Reweighted maps & \ensuremath{-}35 & 0.6 & 12 & 18 & 22 & 32 & 41 & 311 (266--378) \\
\quad No maps & \ensuremath{-}1.3 & 2.5 & 5.9 & 7.6 & 8.9 & 12 & 13 & $>$800 \\
\bottomrule
\end{tabular}

\end{table}

\subsection*{One study, one draw, one atlas}
The validation's intervals resample held-out donors with predictions averaged over five draws of paired cells and one atlas. A study draws its paired cells once. With each draw index scored on its own, and in 2,000 resamples that drew donors and then one draw index per fold (shared by all methods), the saving over paired-only estimation (V1) and the ratio of the paired cells that reference regression and the estimator needed (K2) were recomputed. In ten bootstrap replicates of the atlas's patients, the atlas context was rebuilt and the estimator and the four non-pooled forms of reference regression refitted on the first draw of every fold; a baseline replicate took every patient once through the same construction. Scored one draw at a time, the estimator needed \RvDrawCellsMin{} to \RvDrawCellsMax{} paired cells, the saving over paired-only estimation was at least \RvDrawVOneMin{} to \RvDrawVOneMax, and K2 was \RvDrawKTwoMin{} to \RvDrawKTwoMax{} (resampled, \RvDrawKTwoMed; 95\% interval \RvDrawKTwoLo--\RvDrawKTwoHi, above one in \RvDrawKTwoShare\% of resamples). With the atlas's patients resampled, the estimator needed \RvAtlasCellsMin{} to \RvAtlasCellsMax{} paired cells and K2 was \RvAtlasKTwoMin{} to \RvAtlasKTwoMax{} over the \RvAtlasReps{} replicates (with donors also resampled, \RvAtlasKTwoLo--\RvAtlasKTwoHi; Supplementary Table~\ref{tab:onedraw}). The estimator's small advantage over reference regression given the atlas therefore depends on the draw and on the atlas's patients; its advantage over paired-only estimation does not.

\begin{table}[h]
\caption{\textbf{One study, one draw, one atlas.} Top: each draw index of the validation test scored on its own (the atlas fixed): paired cells the estimator needed to reach the primary target, the saving over the paired-only envelope (V1; a lower bound when paired-only estimation did not reach the target within 800) and the ratio of the paired cells needed by the best form of reference regression to those needed by the estimator (K2); then 2,000 resamples of donors with one draw index per fold, median and 95\% interval. Bottom: the atlas's patients resampled (ten bootstrap replicates; first draw of paired cells), with the replicate that takes every patient once; then the replicates combined with 200 resamples of donors each. The central 95\% of these nested resamples is a sensitivity summary: ten atlas replicates do not establish nominal coverage over new reference atlases.}
\label{tab:onedraw}
\centering\footnotesize
% Generated by extension/synthesis/make_assets.py. Do not edit.
\begin{tabular}{@{}lrrr@{}}
\toprule
 & Paired cells, estimator & Saving over paired-only (V1) & Regression over estimator (K2) \\
\midrule
\multicolumn{4}{@{}l}{\textit{Each draw of paired cells on its own (atlas fixed)}} \\
\quad Draw 1 & 152 & 5.27 & 1.09 \\
\quad Draw 2 & 142 & 5.65 & 1.17 \\
\quad Draw 3 & 154 & 5.19 & 0.95 \\
\quad Draw 4 & 133 & 6.03 & 1.04 \\
\quad Draw 5 & 151 & 5.30 & 1.08 \\
\quad One draw per fold, donors resampled & & 5.53 (4.21--6.68) & 1.04 (0.85--1.45) \\
\addlinespace
\multicolumn{4}{@{}l}{\textit{Atlas patients resampled (first draw)}} \\
\quad Every patient once & 147 & 5.43 & 1.13 \\
\quad Replicate 1 & 166 & 4.83 & 0.99 \\
\quad Replicate 2 & 160 & 5.00 & 1.05 \\
\quad Replicate 3 & 146 & 5.49 & 1.10 \\
\quad Replicate 4 & 170 & 4.71 & 1.00 \\
\quad Replicate 5 & 151 & 5.30 & 1.27 \\
\quad Replicate 6 & 159 & 5.04 & 1.08 \\
\quad Replicate 7 & 141 & 5.68 & 1.18 \\
\quad Replicate 8 & 158 & 5.05 & 1.06 \\
\quad Replicate 9 & 160 & 5.00 & 1.04 \\
\quad Replicate 10 & 152 & 5.28 & 1.13 \\
\quad Replicates and donors resampled & & 5.04 (3.85--6.34) & 1.07 (0.82--1.40) \\
\bottomrule
\end{tabular}

\end{table}

\subsection*{Named within-type relationships}
Eight relationships between a gene program and a protein were fixed from prior knowledge of blood cell states before any was scored (Supplementary Table~\ref{tab:relationships}). A relationship's value in a held-out sample and cell type is the mean, over the program's genes, of their within-type correlation with the protein; a prediction's value is the same mean of the predicted matrix. Endpoints were the recovered fraction of each relationship over held-out units, the share of units in which the prediction had the sign of the held-out value (among units whose two halves agreed in sign) and, per donor, the difference in absolute error between the estimator and the best paired-only method or the best form of reference regression, both chosen per relationship and budget on the pooled error, which favours the comparators. With 100 paired cells, the estimator recovered more of the relationship than reference regression in \RvRelBetterReg{} of the \RvRelN{} relationships and more than paired-only estimation in \RvRelBetterPo; its absolute error was lower than regression's in \RvRelDonorWinReg\% of the comparisons between donors and lower than paired-only estimation's in \RvRelDonorWinPo\%. Every relationship had the expected sign in the held-out cells, two of them weakly: interferon-stimulated genes with CD169 in CD14 monocytes ($r=\RvRelOneTruth$) and cytotoxic genes with CD56 in NK cells ($r=\RvRelSevenTruth$, negative in \RvRelSevenTruthShare\% of units). The second was too weak for 100 paired cells: the estimator's predictions scattered around zero (mean \RvRelSevenAtlasPred) and recovered less than predicting no association (\RvRelSevenAtlasRf\%), and reference regression's did so too, by less (\RvRelSevenRegRf\%).

\begin{table}[h]
\caption{\textbf{Named within-type relationships with 100 paired cells.} Programs: interferon-stimulated (R1: ISG15, IFI6, IFI44L, MX1, OAS1, OAS2, HERC5, XAF1, EPSTI1, IFITM3, STAT2, IRF9), MHC class II (R2, R3: HLA-DRA, HLA-DRB1, HLA-DRB5, HLA-DQA1, HLA-DQB1, HLA-DPA1, HLA-DPB1, HLA-DMA, HLA-DMB), cytotoxic in CD8 T cells (R4, R5: GZMB, GZMH, PRF1, GNLY, NKG7, FGFBP2, CST7, KLRD1, ADGRG1) and in NK cells (R6, R7: GZMB, PRF1, FGFBP2, SPON2, GNLY, NKG7), and LEF1 (R8); expected signs positive except R5 and R7. Units, held-out samples in which the cell type was scored; held-out $r$, the mean over units of the relationship's value; recovered, the fraction of the relationship recovered by the estimator with the atlas (Est.), reference regression with the atlas (Regr.) and paired-only estimation (Paired), each comparator in its best form for the relationship; correct sign, among units whose two halves agreed in sign; donors, the number of donors in which the estimator's absolute error was lower, of all donors with the cell type scored.}
\label{tab:relationships}
\centering\scriptsize
% Generated by extension/synthesis/make_assets.py. Do not edit.
\begin{tabular}{@{}lrrrrrrrrrr@{}}
\toprule
 & & Held-out & \multicolumn{3}{c}{Recovered (\%)} & \multicolumn{3}{c}{Correct sign (\%)} & \multicolumn{2}{c}{Donors, estimator better} \\
\cmidrule(lr){4-6}\cmidrule(lr){7-9}\cmidrule(l){10-11}
Relationship & Units & $r$ & Est. & Regr. & Paired & Est. & Regr. & Paired & than regr. & than paired \\
\midrule
R1\enspace CD14 Mono: interferon, CD169 & 34 & 0.07 & 27 & 21 & 15 & 72 & 79 & 76 & 14/18 & 15/18 \\
R2\enspace CD14 Mono: MHC II, HLA-DR & 34 & 0.47 & 80 & 46 & 19 & 100 & 100 & 93 & 18/18 & 18/18 \\
R3\enspace B: MHC II, HLA-DR & 34 & 0.15 & 81 & 54 & 39 & 100 & 100 & 100 & 17/18 & 17/18 \\
R4\enspace CD8 T: cytotoxic, CD57 & 34 & 0.50 & 65 & 63 & 57 & 99 & 93 & 94 & 12/18 & 12/18 \\
R5\enspace CD8 T: cytotoxic, CD27 & 34 & \ensuremath{-}0.49 & 69 & 66 & 56 & 99 & 100 & 99 & 10/18 & 16/18 \\
R6\enspace NK: cytotoxic, CD16 & 34 & 0.38 & 53 & 37 & 18 & 98 & 95 & 81 & 18/18 & 18/18 \\
R7\enspace NK: cytotoxic, CD56 & 34 & \ensuremath{-}0.05 & \ensuremath{-}81 & \ensuremath{-}20 & 0.0 & 36 & 46 & 0 & 4/18 & 3/18 \\
R8\enspace CD4 T: LEF1, CD45RA & 34 & 0.33 & 55 & 38 & 29 & 100 & 100 & 96 & 18/18 & 18/18 \\
\bottomrule
\end{tabular}

\end{table}

\subsection*{The association test and the truth it defines}
The readout defines a reproducible association by Fisher's $z$ with $P<0.01$ and the same sign in both halves of the held-out cells, down to 20 cells per half. Permuting protein vectors across the cells of each half, which keeps both marginal distributions and removes the pairing (20 permutations per unit and half), gives the test's null distribution under the actual data. The test was calibrated at the thresholds the readout uses: over all \RvFisherUnits{} units, \RvFisherFive\% of null tests gave $P<0.05$ and \RvFisherOne\% gave $P<0.01$, and in every range of cells per half \RvFisherOneMin\% to \RvFisherOneMax\% gave $P<0.01$ (Supplementary Table~\ref{tab:fisher}); the far tail was somewhat heavier (\RvFisherTenth\% at $P<0.001$). No scored unit had fewer than \RvFisherSmallLo{} cells per half, and in the smallest (\RvFisherSmallRange{} cells) \RvFisherSmallOne\% of null tests gave $P<0.01$. Independently permuted halves passed the readout's rule, $P<0.01$ in both with the same sign, for \RvReplNull{} of every 100,000 pairs, against a nominal \RvReplNominal{} and an observed \RvReplObserved\% of pairs, so false associations make up about \RvReplFalseShare\% of those counted.

The readout was also recomputed with the held-out truth centred within the authors' finer cell states (for example naive and memory B cells) instead of the five broad types, with RNA and protein residualized on the log library sizes of both modalities and the 10x lane, and with both; reproducible associations were redefined under each truth, with Fisher's $z$ scaled by $\sqrt{n-3-k}$ for $k$ covariates (Supplementary Table~\ref{tab:truths}). Removing the technical covariates changed little: there were \RvTechReps{} reproducible associations instead of \RvPrimaryReps, and with 100 paired cells the atlas's gain over paired-only estimation was \RvTechOneMinusPo{} points in direction (95\% CI, \RvTechOneMinusPoLo--\RvTechOneMinusPoHi) and \RvTechTwoMinusPo{} in the top pairs. Within finer states, most associations disappeared (\RvFineReps), and the gain fell to \RvFineOneMinusPo{} points in direction (\RvFineOneMinusPoLo--\RvFineOneMinusPoHi) and \RvFineTwoMinusPo{} in the top pairs; the recovered fraction became negative (\RvFineRfAtlas\%, against \RvFineRfPo\% for paired-only estimation and \RvFineRfReg\% for reference regression), because the estimate, which targets the dependence within a broad type, includes the part that lies between the finer states. Much of the within-type dependence that the readout scores thus lies between cell states within the broad types; estimating it within finer states would require centring the paired cells within them.

\begin{table}[h]
\caption{\textbf{The association test under permutation.} Share of Fisher-$z$ tests at or below each $P$ value when protein vectors were permuted across the cells of each half of a held-out unit (20 permutations), by the number of cells per half; tests are gene--protein pairs.}
\label{tab:fisher}
\centering\footnotesize
% Generated by extension/synthesis/make_assets.py. Do not edit.
\begin{tabular}{@{}lrrrr@{}}
\toprule
 & & \multicolumn{3}{c}{Share of null tests at or below (\%)} \\
\cmidrule(l){3-5}
Cells per half & Tests & 5\% & 1\% & 0.1\% \\
\midrule
40--79 & 888,000 & 4.54 & 1.07 & 0.18 \\
80--159 & 17,760,000 & 4.71 & 1.10 & 0.18 \\
160 or more & 132,312,000 & 4.92 & 1.06 & 0.14 \\
All & 150,960,000 & 4.90 & 1.07 & 0.15 \\
\bottomrule
\end{tabular}

\end{table}

\begin{table}[h]
\caption{\textbf{The readout under other held-out truths.} The atlas estimate minus paired-only estimation (best of eight methods at each budget) and minus reference regression (best of six forms), with 100 paired cells and 95\% intervals from resampling donors, in the share of reproducible associations whose direction was predicted (E1), the share of reproducible associations of the predicted direction among the 100 largest predictions (E2) and the recovered fraction. Truths: centred within the five broad types (the published readout), within the authors' finer cell states, within broad types after residualizing on log library sizes and the 10x lane, and both.}
\label{tab:truths}
\centering\scriptsize
% Generated by extension/synthesis/make_assets.py. Do not edit.
\begin{tabular}{@{}lrrlrrr@{}}
\toprule
 & & & & \multicolumn{3}{c}{Estimator minus comparator (points)} \\
\cmidrule(l){5-7}
Held-out truth & Units & Associations & Comparator & Direction (E1) & Top 100 (E2) & Recovered \\
\midrule
Broad types (published) & 170 & 120,649 & Paired only & 11 (10 to 12) & 19 (18 to 20) & 19 (17 to 20) \\
 &  &  & Regression & \ensuremath{-}0.1 (\ensuremath{-}0.5 to 0.2) & \ensuremath{-}2.8 (\ensuremath{-}3.6 to \ensuremath{-}1.9) & 0.1 (\ensuremath{-}1.3 to 1.6) \\
\addlinespace
Finer states & 170 & 31,509 & Paired only & 2.8 (2.1 to 3.4) & 6.5 (6.1 to 7.0) & \ensuremath{-}71 (\ensuremath{-}87 to \ensuremath{-}58) \\
 &  &  & Regression & \ensuremath{-}2.0 (\ensuremath{-}2.7 to \ensuremath{-}1.3) & \ensuremath{-}1.9 (\ensuremath{-}2.2 to \ensuremath{-}1.6) & \ensuremath{-}70 (\ensuremath{-}85 to \ensuremath{-}57) \\
\addlinespace
Technical covariates & 170 & 98,739 & Paired only & 10 (9.6 to 11) & 18 (17 to 19) & 13 (11 to 16) \\
 &  &  & Regression & 0.1 (\ensuremath{-}0.3 to 0.5) & \ensuremath{-}1.7 (\ensuremath{-}2.3 to \ensuremath{-}1.1) & \ensuremath{-}3.9 (\ensuremath{-}5.8 to \ensuremath{-}2.0) \\
\addlinespace
Both & 170 & 26,750 & Paired only & 3.2 (2.6 to 3.6) & 5.6 (5.1 to 6.1) & \ensuremath{-}89 (\ensuremath{-}107 to \ensuremath{-}73) \\
 &  &  & Regression & \ensuremath{-}0.6 (\ensuremath{-}1.2 to \ensuremath{-}0.0) & \ensuremath{-}1.0 (\ensuremath{-}1.2 to \ensuremath{-}0.7) & \ensuremath{-}87 (\ensuremath{-}104 to \ensuremath{-}71) \\
\bottomrule
\end{tabular}

\end{table}

\subsection*{The law with data sets as units}
The prospective test's rank correlation was tested by permuting problems, but panels of 100, 200 and 400 features of one data set are nested, and the ten problems with a measurable saving came from four data sets of two studies (mouse lymphoid CITE-seq with two antibody panels; fly visual neurons in two connectomes). The law's error was therefore summarized per data set and without each data set in turn, the rank correlation was computed between data-set means with an exact permutation test over data sets and between problems with permutations only within data sets, and the slope of observed on predicted saving was estimated with a fixed effect per data set (Supplementary Table~\ref{tab:lawunits}). At half of the dependence, the law overpredicted in all \RvLawHalfSets{} data sets, with median errors of \RvLawHalfLooMin\% to \RvLawHalfLooMax\% when any one was left out (95\% CI of the median error \RvLawHalfErrLo--\RvLawHalfErrHi\% resampling data sets and \RvLawHalfErrStLo--\RvLawHalfErrStHi\% resampling studies). The data sets' averages were ranked in order ($\rho=\RvLawHalfMeansRho$, exact $P=\RvLawHalfMeansP$, the smallest attainable with four), the problems were ranked within data sets ($P=\RvLawHalfWithinP$ over all \RvLawHalfWithinPerm{} permutations within data sets), and the slope with data-set effects was \RvLawHalfSlope{} (\RvLawHalfSlopeLo--\RvLawHalfSlopeHi). At 30\% of the dependence the law overpredicted in all \RvLawThreeSets{} data sets with a measurable saving (median error \RvLawThreeErr\%); at 70\%, only the two fly data sets had one (\RvLawSevenErr\%). Of the \RvLawHalfNonN{} problems without a measurable saving at half of the dependence, \RvLawHalfBelow{} had both curves above the target with 25 paired cells, \RvLawHalfNotReached{} had a curve that did not reach it within the budgets, and \RvLawHalfOneAbove{} had one curve starting above it.

\begin{table}[h]
\caption{\textbf{The law with data sets as units.} Per data set with a measurable saving: problems, median absolute error of the law ($2^{|\log_2(\text{observed}/\text{law})|}-1$), median ratio of the observed to the predicted saving, and the median error over the other data sets. All: the median error with 95\% intervals from resampling data sets and, in the same parentheses, studies. Then the rank correlation of data-set means with its exact permutation $P$ value, the problem-level rank correlation with permutations within data sets, the slope of $\log$ observed on $\log$ predicted saving with a fixed effect per data set (95\% interval from resampling data sets), and the problems without a measurable saving by reason.}
\label{tab:lawunits}
\centering\footnotesize
% Generated by extension/synthesis/make_assets.py. Do not edit.
\begin{tabular}{@{}lrrrr@{}}
\toprule
Data set & Problems & Median error (\%) & Observed / law & Without this data set (\%) \\
\midrule
\multicolumn{5}{@{}l}{\textit{Half of the dependence recovered (primary)}} \\
\quad Fly visual neurons, BANC & 2 & 10 & 0.91 & 25 \\
\quad Fly visual neurons, FlyWire & 3 & 18 & 0.85 & 26 \\
\quad Mouse lymphoid organs, 111 proteins & 3 & 26 & 0.80 & 18 \\
\quad Mouse lymphoid organs, 206 proteins & 2 & 35 & 0.74 & 17 \\
\quad All & 10 & 21 (11--29; studies 13--26) & & \\
\multicolumn{5}{@{}p{.95\textwidth}}{\quad rank correlation of data-set means 1.00 (exact $P=0.042$, 4 data sets); problem-level $\rho=0.98$, permutations within data sets $P=0.007$ (144); slope with data-set effects 1.08 (0.99--1.26).} \\
\multicolumn{5}{@{}p{.95\textwidth}}{\quad Not measurable: 10 a curve does not reach the target within the budgets; 8 target below both curves at the smallest budget; 1 one curve starts above the target.} \\
\midrule
\multicolumn{5}{@{}l}{\textit{30\% recovered}} \\
\quad Human cortex Patch-seq & 2 & 144 & 0.41 & 61 \\
\quad Mouse lymphoid organs, 111 proteins & 3 & 68 & 0.60 & 81 \\
\quad Mouse lymphoid organs, 206 proteins & 3 & 54 & 0.65 & 99 \\
\quad All & 8 & 74 (54--144; studies 61--144) & & \\
\multicolumn{5}{@{}p{.95\textwidth}}{\quad rank correlation of data-set means 1.00 (exact $P=0.167$, 3 data sets); problem-level $\rho=0.95$, permutations within data sets $P=0.028$ (72); slope with data-set effects 0.98 (-2.39--1.22).} \\
\multicolumn{5}{@{}p{.95\textwidth}}{\quad Not measurable: 9 target below both curves at the smallest budget; 7 a curve does not reach the target within the budgets; 5 one curve starts above the target.} \\
\midrule
\multicolumn{5}{@{}l}{\textit{70\% recovered}} \\
\quad Fly visual neurons, BANC & 2 & 2 & 1.02 & 5 \\
\quad Fly visual neurons, FlyWire & 3 & 5 & 0.95 & 2 \\
\quad All & 5 & 4 (2--5; studies 4--4) & & \\
\multicolumn{5}{@{}p{.95\textwidth}}{\quad problem-level $\rho=1.00$, permutations within data sets $P=0.083$ (12); slope with data-set effects 1.02 (0.99--1.15).} \\
\multicolumn{5}{@{}p{.95\textwidth}}{\quad Not measurable: 15 a curve does not reach the target within the budgets; 9 one curve starts above the target.} \\
\bottomrule
\end{tabular}

\end{table}

\subsection*{A pilot that sets the budget}
For each of the 29 problems of the prospective test, target accuracy (a recovered fraction of 0.3, 0.5 or 0.7) and pilot size (50, 100 or 200 paired cells, at most a fifth of the reservoir), 100 replicates drew a random pilot from the reservoir, estimated each block's signal and noise from it, set the total budget to the paired cells that Proposition~\ref{prop:law} gives for block shrinkage at the target (never fewer than the pilot), drew the remaining cells so that the pilot counted towards the total, fitted block James--Stein to all of them and scored the recovered fraction against the truth cells. A replicate whose pilot showed no positive signal, or whose budget exceeded the reservoir, used the whole reservoir and was flagged. The same fit at the budget chosen from all reservoir cells served as reference (Supplementary Table~\ref{tab:pilotdesign}). At half of the dependence, pilots of 50, 100 and 200 cells chose budgets that reached the target in \RvPilotHalfFiftyReached\%, \RvPilotHalfHundredReached\% and \RvPilotHalfTwoHundredReached\% of replicates, and reached at least the target less 5 percentage points in \RvPilotHalfFiftyWithin\%, \RvPilotHalfHundredWithin\% and \RvPilotHalfTwoHundredWithin\%. Their median totals were \RvPilotHalfFiftyBudget, \RvPilotHalfHundredBudget{} and \RvPilotHalfTwoHundredBudget{} paired cells: from 100 cells on, the pilot alone usually exceeded the budget that the law gave. The budget from all reservoir cells was a median of \RvPilotHalfResRatio{} times the paired cells that block shrinkage needed, and it reached the target in \RvPilotThreeResReached\%, \RvPilotHalfResReached\% and \RvPilotSevenResReached\% of replicates at 30\%, 50\% and 70\%: the law assumes the best factor in every block, and the positive-part rule needs more cells than that factor would. At 70\%, \RvPilotSevenFiftyFlagged\% to \RvPilotSevenTwoHundredFlagged\% of replicates were flagged. By data set, pilots fell more than 5 points short of the target in more than half of the replicates only in human cortical Patch-seq and in the chromatin of fetal cortex at 50\%, and in these and the chromatin of adult mouse cortex at 70\%.

\begin{table}[h]
\caption{\textbf{A pilot that sets the budget.} For each target and pilot size, over the 29 problems of the prospective test (each weighted alike; problems whose reservoir was too small for a pilot size used the largest pilot allowed, a fifth of the reservoir): the median over problems of each problem's median achieved recovered fraction, with the interquartile range over problems; the share of replicates reaching the target, and the share reaching at least the target less 5 percentage points (overshooting included); the median total budget, pilot included; its median ratio to the paired cells that the observed curve needed (problems whose curve crossed the target); and the share of replicates flagged (no positive signal in the pilot, or a budget beyond the reservoir, so that the whole reservoir was used). All reservoir cells: the budget chosen with the block moments of every reservoir cell.}
\label{tab:pilotdesign}
\centering\scriptsize
% Generated by extension/synthesis/make_assets.py. Do not edit.
\begin{tabular}{@{}llrrrrrrr@{}}
\toprule
 & & & Achieved & Reached & Reached target & & Budget / & Flagged \\
Target & Pilot & Problems & (\%; IQR) & (\%) & $-$5 points (\%) & Budget & needed & (\%) \\
\midrule
30\% & 50 cells & 29 & 41 (34--90) & 73 & 78 & 50 & 0.99 & 0 \\
 & 100 cells & 29 & 41 (35--94) & 77 & 82 & 100 & 1.02 & 1 \\
 & 200 cells & 29 & 48 (34--97) & 83 & 87 & 200 & 1.16 & 1 \\
 & All reservoir cells & 29 & 23 (12--26) & 34 & 45 & 32 & 0.57 & 0 \\
\addlinespace
50\% & 50 cells & 29 & 54 (48--90) & 64 & 72 & 84 & 1.16 & 5 \\
 & 100 cells & 29 & 53 (48--94) & 65 & 75 & 100 & 1.12 & 8 \\
 & 200 cells & 29 & 53 (48--97) & 65 & 77 & 200 & 0.87 & 10 \\
 & All reservoir cells & 29 & 46 (40--48) & 26 & 50 & 86 & 0.77 & 0 \\
\addlinespace
70\% & 50 cells & 29 & 68 (60--90) & 46 & 56 & 181 & 0.82 & 22 \\
 & 100 cells & 29 & 69 (61--94) & 48 & 59 & 181 & 0.85 & 24 \\
 & 200 cells & 29 & 70 (61--97) & 56 & 61 & 200 & 0.99 & 25 \\
 & All reservoir cells & 29 & 67 (61--69) & 23 & 49 & 181 & 0.82 & 0 \\
\bottomrule
\end{tabular}

\end{table}

\subsection*{Savings at fixed accuracy}
The savings summary uses each data set's own target and counts curves that never reach it at the largest budget. At fixed recovered fractions from 5\% to 40\%, the paired cells of the estimator and of paired-only estimation were read from the published curves both with and without the running maximum, and each comparison was classed as estimable (both curves cross the level between two budgets), unresolved above (a curve never reaches it) or unresolved below (both start above it); estimable ratios have 95\% intervals from the tests' resamples of held-out units (Supplementary Tables~\ref{tab:fixed} and~\ref{tab:fixeddetail}). Read from the raw curves, the saving could be measured in \RvFixedTenN, \RvFixedTwentyN{} and \RvFixedThirtyN{} of the \RvFixedSets{} data sets at 10\%, 20\% and 30\% recovery (geometric means \RvFixedTenGm, \RvFixedTwentyGm{} and \RvFixedThirtyGm), and the running maximum changed \RvFixedSameRunningMax{} of the counts. At 20\%, the other comparisons were unresolved because paired-only estimation did not reach the level within 800 paired cells (\RvFixedTwentyPoAbove{} data sets), because neither method did (\RvFixedTwentyBothAbove) or because the estimator had already exceeded it with 25 paired cells (\RvFixedTwentyEstBelow). The summary's savings, read at each data set's own target and counting curves that never reach it, therefore rest on comparisons that fixed levels mostly leave unresolved.

\begin{table}[h]
\caption{\textbf{Savings at fixed accuracy.} Over the nine benchmark data sets and the external test, at each recovered fraction: comparisons that could be made (both the estimator and the paired-only envelope cross the level between two budgets), the geometric mean of their savings, and comparisons left unresolved because a curve never reached the level (above) or one or both curves started above it with 25 paired cells (below), read from the raw curves and from their running maximum.}
\label{tab:fixed}
\centering\footnotesize
% Generated by extension/synthesis/make_assets.py. Do not edit.
\begin{tabular}{@{}lrrrrrrrr@{}}
\toprule
 & \multicolumn{4}{c}{Raw curves} & \multicolumn{4}{c}{Running maximum} \\
\cmidrule(lr){2-5}\cmidrule(l){6-9}
Recovered & Estimable & Geometric mean & Above & Below & Estimable & Geometric mean & Above & Below \\
\midrule
5\% & 3 & 4.29 & 4 & 3 & 3 & 4.29 & 4 & 3 \\
10\% & 2 & 1.74 & 6 & 2 & 2 & 1.74 & 6 & 2 \\
15\% & 2 & 2.07 & 6 & 2 & 2 & 2.07 & 6 & 2 \\
20\% & 2 & 2.46 & 6 & 2 & 2 & 2.46 & 6 & 2 \\
30\% & 1 & 4.88 & 7 & 2 & 1 & 4.88 & 7 & 2 \\
40\% & 2 & 3.39 & 8 & 0 & 2 & 3.39 & 8 & 0 \\
\bottomrule
\end{tabular}

\end{table}

\begin{table}[h]
\caption{\textbf{Savings at fixed accuracy per data set.} Ratio of the paired cells needed by paired-only estimation to those needed by the estimator at 10\%, 20\% and 30\% of the dependence recovered, read from the raw curves, with 95\% intervals from the tests' resamples of held-out units; otherwise the reason the comparison is unresolved.}
\label{tab:fixeddetail}
\centering\footnotesize
% Generated by extension/synthesis/make_assets.py. Do not edit.
\begin{tabular}{@{}lrrr@{}}
\toprule
Data set & 10\% recovered & 20\% recovered & 30\% recovered \\
\midrule
External test (OverCITE-seq) & 2.62 (2.06--3.27) & 4.35 (3.21--5.19) & 4.88 (3.22--6.03) \\
Blood (Hao) & paired only not reached & paired only not reached & paired only not reached \\
Blood (Stephenson) & not reached & not reached & not reached \\
Bone marrow CITE-seq & paired only not reached & paired only not reached & paired only not reached \\
Colon & paired only not reached & paired only not reached & not reached \\
Bone marrow RNA--chromatin & paired only not reached & paired only not reached & not reached \\
Patch-seq, visual cortex & 1.15 (0.97--1.33) & 1.40 (1.17--1.59) & paired only not reached \\
Patch-seq, motor cortex & not reached & not reached & not reached \\
Fly, female & both above at 25 & estimator above at 25 & estimator above at 25 \\
Fly, male & both above at 25 & estimator above at 25 & estimator above at 25 \\
\bottomrule
\end{tabular}

\end{table}
